\documentclass[10pt]{article}
\usepackage[T1]{fontenc}\usepackage[utf8]{inputenc}\usepackage{textcomp}
\usepackage[scaled]{helvet}\renewcommand{\familydefault}{\sfdefault}
\DeclareUnicodeCharacter{03B5}{\ensuremath{\varepsilon}}\DeclareUnicodeCharacter{03B2}{\ensuremath{\beta}}
\DeclareUnicodeCharacter{03B7}{\ensuremath{\eta}}\DeclareUnicodeCharacter{03B1}{\ensuremath{\alpha}}
\DeclareUnicodeCharacter{03C7}{\ensuremath{\chi}}\DeclareUnicodeCharacter{2212}{\ensuremath{-}}
\DeclareUnicodeCharacter{2265}{\ensuremath{\geq}}\DeclareUnicodeCharacter{2264}{\ensuremath{\leq}}
\DeclareUnicodeCharacter{2192}{\ensuremath{\rightarrow}}
\usepackage[a4paper,margin=20mm]{geometry}
\usepackage{graphicx}
\pagestyle{empty}\setlength{\parindent}{0pt}\setlength{\parskip}{3pt}
\begin{document}
{\large\bfseries Genome-wide association and haplotype analysis identify candidate genes for grain nutritional quality in wild barley (\textit{Hordeum spontaneum})\par}\vspace{3mm}
Theoretical and Applied Genetics\par
[AUTHOR NAMES]\par
Corresponding author: [CORRESPONDING AUTHOR], [AFFILIATION: institution, department, city, country]; e-mail: [E-MAIL]\par\vspace{10mm}
{\large\bfseries Online Resource 9\par}\vspace{3mm}
Haplotypes of the four presented genes, as plotted by crosshap v1.4.0 at the settings used in the analysis (SNPs within the gene ± 1 kb): \textbf{a} \textit{GPAT6} (HORVU.MOREX.r3.3HG0301300), fiber; \textbf{b} \textit{GH17} (HORVU.MOREX.r3.5HG0487060), fiber; \textbf{c} \textit{PHT4;3} (HORVU.MOREX.r3.3HG0301710), starch (ε = 0.6, MGmin = 2); the GDSL esterase/lipase gene HORVU.MOREX.r3.7HG0729030, \textbf{d} fiber and \textbf{e} starch (ε = 0.9, MGmin = 3; the same haplotypes for both traits)

In each plot: top, the number of accessions in each haplotype group (A, B, …); center, the marker-group (MG) alleles that define each haplotype (open circle, reference allele; filled circle, alternate allele); left, the genotype counts of the SNPs of each marker group (reference, missing, heterozygous, alternate); right, the absolute difference in the trait between alternate- and reference-allele carriers at each SNP of the marker group (Pheno association), shaded by alternate-allele frequency; bottom, the genotype BLUPs of the trait (\% of grain) in each haplotype group. Tables: for each marker group, its number of SNPs (nSNP), phenotypic difference (phenodiff), mean r² (meanR2) and alternate-allele frequency (AltAF); for each haplotype group, its mean BLUP (Pheno) and number of accessions (nTotal). Accessions not assigned to a haplotype group are not shown. The haplotype groups are those of Figs. 3 and 4; their tests are in Online Resource 6 and Fig. 4\par\vspace{6mm}
{\small Page 2: a, GPAT6 (fiber)\\ Page 3: b, GH17 (fiber)\\ Page 4: c, PHT4;3 (starch)\\ Page 5: d, GDSL esterase/lipase (fiber)\\ Page 6: e, GDSL esterase/lipase (starch)}
\end{document}
